\chapter{Штучний інтелект та інтелектуальні сервіси}

\section{Архітектура локального ШІ-інференсу}
У сучасних автономних інформаційних системах (Air-Gapped) використання зовнішніх хмарних моделей ШІ (наприклад, OpenAI або Anthropic) є неприпустимим через вимоги конфіденційності та безпеки. Підсистема «ERP/1: Істотність» базується на використанні локальних компактних моделей ШІ (класу Llama та Phi) на споживчих графічних процесорах (Consumer GPU), що мають лімітований об'єм оперативної пам'яті (до 16 ГБ VRAM).

Архітектура побудована на базі сервісу інференсу \texttt{llama.cpp} над операційною системою Ubuntu 24.04 з підтримкою драйверів NVIDIA CUDA 12.x. Інтеграція з основним Erlang/OTP середовищем здійснюється через асинхронні REST API / WebSocket запити за допомогою легковагого HTTP-клієнта.

\section{Математична модель розрахунку об'єму відеопам'яті (VRAM)}
Для забезпечення стабільної паралельної роботи користувачів без ризику виникнення помилок Out-Of-Memory (OOM), здійснюється чітке планування та резервування VRAM. Загальний об'єм відеопам'яті ($V_{\text{VRAM}}$) розподіляється на три ключові компоненти:
\begin{equation}
    V_{\text{VRAM}} \ge V_{\text{model}} + V_{\text{KV}} + V_{\text{workspace}}
\end{equation}

Де:
\begin{itemize}
    \item $V_{\text{model}}$ --- об'єм пам'яті для зберігання статичних ваг моделі з урахуванням квантування;
    \item $V_{\text{KV}}$ --- динамічний буфер пам'яті для зберігання контексту користувачів (Key-Value Cache);
    \item $V_{\text{workspace}}$ --- системний робочий буфер CUDA для виконання обчислень (зазвичай резервується 512--1024 МБ).
\end{itemize}

\subsection{Обчислення для моделі Llama-3-8B-Instruct}
При квантуванні Q4\_K\_M статична вага моделі становить $V_{\text{model}} \approx 4.8\text{ ГБ}$. Розмір KV-кешу для одного запиту з контекстом $N_{\text{ctx}} = 8192$ токенів розраховується за формулою:
\begin{equation}
    V_{\text{KV\_single}} = 2 \times N_{\text{layers}} \times N_{\text{heads}} \times d_{\text{head}} \times N_{\text{ctx}} \times N_{\text{bytes\_per\_elem}}
\end{equation}

Для Llama-3-8B ($N_{\text{layers}} = 32$, $N_{\text{heads}} = 8$, $d_{\text{head}} = 128$, квантований KV-кеш у форматі 16-bit FP):
\begin{equation}
    V_{\text{KV\_single}} = 2 \times 32 \times 8 \times 128 \times 8192 \times 2 \approx 1.07\text{ ГБ}
\end{equation}

На споживчій відеокарті NVIDIA RTX 4070 Ti (12 ГБ VRAM) при $V_{\text{workspace}} = 1.0\text{ ГБ}$:
\begin{equation}
    V_{\text{KV\_total}} \le 12.0 - 4.8 - 1.0 = 6.2\text{ ГБ}
\end{equation}
Максимальна кількість паралельних сесій: $M = \lfloor 6.2 / 1.07 \rfloor = 5$ користувачів.

\subsection{Обчислення для моделі Phi-3-medium-14B}
При квантуванні Q4\_K\_M статична вага моделі становить $V_{\text{model}} \approx 8.5\text{ ГБ}$. При контексті $N_{\text{ctx}} = 4096$:
\begin{equation}
    V_{\text{KV\_single}} = 2 \times 40 \times 10 \times 96 \times 4096 \times 2 \approx 0.625\text{ ГБ}
\end{equation}

На відеокарті NVIDIA RTX 4070 Ti Super (16 ГБ VRAM) при $V_{\text{workspace}} = 1.0\text{ ГБ}$:
\begin{equation}
    V_{\text{KV\_total}} \le 16.0 - 8.5 - 1.0 = 6.5\text{ ГБ}
\end{equation}
Максимальна кількість паралельних сесій: $M = \lfloor 6.5 / 0.625 \rfloor = 10$ користувачів.

\section{Механізми оптимізації інференсу}
Для ефективної утилізації обчислювальних ядер GPU застосовуються такі технології:
\begin{enumerate}
    \item \textbf{Continuous Batching (Динамічне пакетування)}: об'єднання запитів на рівні окремих ітерацій генерації токенів, що усуває простої під час виконання фаз prefill (обробка промпту) та decode (генерація наступних токенів).
    \item \textbf{PagedAttention}: розбиття KV-кешу на невеликі сторінки фіксованого розміру (наприклад, по 16 токенів) та їх динамічний розподіл через віртуальну таблицю вказівок. Це повністю усуває фрагментацію пам'яті та дозволяє спільно використовувати префікси промптів різними користувачами.
\end{enumerate}
